\scanpage{35}

\diagram{\begin{tikzcd}[ampersand replacement=\&,column sep=small]
\& X\arrow[dl,dashed,"\varphi"']\arrow[d] \& Y\arrow[dl]\\
S \& T\arrow[l,"f"] \& T\times_S T\arrow[l,bend right=15,"\mathrm{pr}_1"']\arrow[l,bend left=15,"\mathrm{pr}_2"]\arrow[ll,bend left=37,"g"]
\end{tikzcd}}

\[
f_*(X/T)=\underline{\mathrm{Hom}}_{T/S}(T/S,X/S)
\]

\[
f_*(X/T)=\underline{\mathrm{Hom}}_{S}(T,X)\times_{\underline{\mathrm{Hom}}_{\underline S}(T,T)}S
\]

\[
\underline{\mathrm{Hom}}_{T/S}(X/T,Y/T)=f_*(\underline{\mathrm{Hom}}_T(X,Y))
\]

\[
\underline{\mathrm{Hom}}_S(X,Y)=\varphi_*(X\times Y/X)
\]

\[
\underline{\mathrm{Hom}}_S(X,Y)=\underline{\mathrm{Hom}}_{S/S}(X/S,Y/S)
\]

Si $X$ est muni d'une donnée de descente : $S$, et si $f_*(X)$ existe\inserted{[et si on est \unsure{sous} les conditions d'\unsure{unicité|université|immersivité}....]}, alors l'objet descendu $X_0$ existe aussi, et on aura

\[
X_0=f_*(X)\times_{g_*(X\times_S T)}f_*(X)
\]

Contre-exemple de \uline{Nagata} : schéma algébrique simple \inserted{\unsure{absolt non-sing|abstr. non-sing}} $X$, sur un corps alg. clos $k$, avec groupe $G\simeq\mathbb{Z}/2$ opérant sans point fixe, tel que $X/G$ n'existe pas. Soit $\deleted{\unsure{K|k'}}L$ une extension de $k$ sur laquelle $G$ opère \deleted{\unsure{sans}}\inserted{\unsure{fidèlmt|fidèlem|fidèlement}} , \deleted{pt fixe} $K=L^G$, faisons opérer $G$ \uline{\unsure{diagonalemt|diagonalement}} sur $X\otimes_{\unsure{K|k}}L\,/\,L$, on a donc sur $X\otimes_k L/L$ une donnée de \uline{descente galoisienne} pour $L/K$. Je dis que \uline{le quotient n'existe pas}, i.e. \uline{la descente n'est pas possible}. En effet, \deleted{si le quotient existait, il existerait aussi par extension des scalaires \inserted{de $K$} à $\overline L$, or}

\deleted{\[
(X\otimes_kL)\otimes_K\unsure{L|\overline L}=X\otimes_k(L\otimes_KL)=X\otimes_kL
\]}

% archive sheet number: 34
